\documentclass[12pt]{article}
\usepackage[a4paper,top=3cm,right=2cm,left=2cm,bottom=2.5cm]{geometry}
\usepackage[latin1]{inputenc}
\usepackage[T1]{fontenc}
\usepackage{newtxtext}
\usepackage[ttscale=0.875]{libertine}
\input{math.tex}
\usepackage[libertine,cmintegrals,cmbraces]{newtxmath}
\usepackage[round]{natbib}
\usepackage[american]{babel}
\usepackage[small]{titlesec}
\usepackage{amsmath}
\usepackage{iexec}
\usepackage{listings}
%% \usepackage{subscript}
%% \usepackage{enumerate}
%% \usepackage{booktabs}
%% \usepackage{tabu}
%% \usepackage{setspace}
%% \usepackage{comment}
\usepackage[defaultlines=4,all]{nowidow}
\usepackage{needspace}
\usepackage{graphicx}
\usepackage{mathrsfs}
\usepackage{ifpdf}
\usepackage{upgreek}
\usepackage[pdftex,colorlinks=true,%
   urlcolor=blue,
   linkcolor=blue,
   citecolor=blue,
   filecolor=blue]{hyperref}
\input{mylistings.tex}
\input{ira.tex}
\renewcommand{\ttdefault}{lmtt}% default cmtt
\newcommand{\texttk}[1]{\texttt{\bfseries #1}}



\title{Documentation for \texttt{nmth.chpl}}

\author{Nelson Lu�s Dias (\texttt{nelsonluisdias@gmail.com})}

\date{\today}

\begin{document}

\maketitle

%% \section*{Notation}

%% In this document,
%% \begin{align}
%%   \Exv\{X\} &= \int_{\mathbb{R}} xf_X(x)\,\md{x}, \\
%%   \Var\{X\} &= \int_{\mathbb{F}} \left(x -
%%   \Exv\{x\}\right)^2f_X(x)\,\md{x}
%% \end{align}
%% (where $f_X(x)$ is the probability density function of $X$) are
%% population parameters, whereas
%% \begin{align}
%%   \Exv\{x\} = \overline{x} = \frac{1}{n}\sum_{i=1}^n x_i, \\
%%   \Var\{x\} = \frac{1}{n}\sum_{i=1}^n \left(x_i-\overline{x}\right)^2
%% \end{align}
%% are sample statistics.



\section{\texttt{logmod}: log-modulus function}

\subsection{Use}
\iexec[quiet,trace,stdout=xtract.txt]{grep 'proc logmod' -A2 nmth.chpl}

\lstinputlisting[style=lstchapelf, numbers=none]{xtract.txt}

\subsection{Theory}

The \emph{log-modulus} function was introduced by
\cite{john1980alternative}. It is asymptotically equal to $x$ for
small values of its argument, and asymptoticall equal to $\ln(\cdot)$
for large positive values of its argument:
\begin{equation}
\logmod(x) \equiv \sign(x)\ln(|x| + 1).\label{eq:def-logmod}
\end{equation}
Note that, asymptotically,
\begin{equation}
  \logmod(x) \sim \begin{cases}
    x, & x \sim 0, \\
    \sign(x)\ln(|x|), & |x| \gg 1,
  \end{cases}
\end{equation}
and that $\logmod(x)$ is odd.




\section{\texttt{expmod}: exponential-modulus function}

\subsection{Use}
\iexec[quiet,trace,stdout=xtract.txt]{grep 'proc expmod' -A2 nmth.chpl}

\lstinputlisting[style=lstchapelf, numbers=none]{xtract.txt}

\subsection{Theory}


We often need the inverse of $\logmod$:
\begin{align*}
  y &= \sign(x)\ln(|x| + 1),  \\
  \sign(y)y &= \ln(|x| + 1),  \\
  \intertext{(because $\logmod$ is odd)}
  \exp\left(\sign(y)y\right) &= |x| + 1, \\
  \exp\left(|y|\right) &= |x| + 1, \\
  |x| &= \exp\left(|y|\right) - 1, \\
  x &= \sign(y)\left[ \exp\left(|y|\right) - 1\right],
\end{align*}
whence
\begin{equation}
\expmod(x) \equiv \sign(x)\left[ \exp\left(|x|\right) - 1\right]\label{eq:def-expmod}
\end{equation}
is the \emph{exponential-modulus} function.

\bibliography{abbrevs,all}
\bibliographystyle{apalike}


\end{document}
